% Options for packages loaded elsewhere
\PassOptionsToPackage{unicode}{hyperref}
\PassOptionsToPackage{hyphens}{url}
\documentclass[
]{article}
\usepackage{xcolor}
\usepackage{amsmath,amssymb}
\setcounter{secnumdepth}{-\maxdimen} % remove section numbering
\usepackage{iftex}
\ifPDFTeX
  \usepackage[T1]{fontenc}
  \usepackage[utf8]{inputenc}
  \usepackage{textcomp} % provide euro and other symbols
\else % if luatex or xetex
  \usepackage{unicode-math} % this also loads fontspec
  \defaultfontfeatures{Scale=MatchLowercase}
  \defaultfontfeatures[\rmfamily]{Ligatures=TeX,Scale=1}
\fi
\usepackage{lmodern}
\ifPDFTeX\else
  % xetex/luatex font selection
\fi
% Use upquote if available, for straight quotes in verbatim environments
\IfFileExists{upquote.sty}{\usepackage{upquote}}{}
\IfFileExists{microtype.sty}{% use microtype if available
  \usepackage[]{microtype}
  \UseMicrotypeSet[protrusion]{basicmath} % disable protrusion for tt fonts
}{}
\makeatletter
\@ifundefined{KOMAClassName}{% if non-KOMA class
  \IfFileExists{parskip.sty}{%
    \usepackage{parskip}
  }{% else
    \setlength{\parindent}{0pt}
    \setlength{\parskip}{6pt plus 2pt minus 1pt}}
}{% if KOMA class
  \KOMAoptions{parskip=half}}
\makeatother
\usepackage{longtable,booktabs,array}
\newcounter{none} % for unnumbered tables
\usepackage{calc} % for calculating minipage widths
% Correct order of tables after \paragraph or \subparagraph
\usepackage{etoolbox}
\makeatletter
\patchcmd\longtable{\par}{\if@noskipsec\mbox{}\fi\par}{}{}
\makeatother
% Allow footnotes in longtable head/foot
\IfFileExists{footnotehyper.sty}{\usepackage{footnotehyper}}{\usepackage{footnote}}
\makesavenoteenv{longtable}
\setlength{\emergencystretch}{3em} % prevent overfull lines
\providecommand{\tightlist}{%
  \setlength{\itemsep}{0pt}\setlength{\parskip}{0pt}}
\usepackage{bookmark}
\IfFileExists{xurl.sty}{\usepackage{xurl}}{} % add URL line breaks if available
\urlstyle{same}
\hypersetup{
  hidelinks,
  pdfcreator={LaTeX via pandoc}}

\author{}
\date{}

\begin{document}

\section{Power Analysis --- Overview}\label{power-analysis-overview}

\subsection{Reviewer Comment}\label{reviewer-comment}

\begin{quote}
The manuscript does not include a power analysis or simulation study
demonstrating that the available sample sizes are adequate for
recovering variance components of the magnitude of interest, and this
omission is consequential in at least two respects. First, the
gender-stratified WLS subsamples are asked to support estimation of
three multilevel models each containing up to nine PGIs, their pairwise
interactions with one another, and their interactions with ascribed
characteristics, a parameter burden that is difficult to reconcile with
a family-level sample of this size; the standard errors of 0.06 on the
individual IO estimates for WLS men, large relative to the point
estimates themselves, are consistent with this concern. Second, and more
fundamentally, the within-family design mechanically reduces the PGI
signal available for estimation---within-family PGI variance is smaller
than between-family variance while the absolute magnitude of measurement
error remains constant---which implies that greater sample sizes are
required to achieve equivalent statistical precision compared to a
between-family design; the manuscript acknowledges this as a measurement
error issue but does not discuss its implications for power.
\end{quote}

\begin{center}\rule{0.5\linewidth}{0.5pt}\end{center}

\subsection{Power Analysis}\label{power-analysis}

\subsubsection{Estimands}\label{estimands}

Two quantities are assessed:

\begin{enumerate}
\def\labelenumi{\arabic{enumi}.}
\item
  \textbf{\texttt{delta\_within\ =\ (emptyind\ -\ condind)\ /\ totalvar}}
  --- within-family variance reduction attributable to PGIs: (m0 → m1).
  Explicitly addresses the reviewer's concern, but is not related to the
  full model with all interactions, which is the real threat to power.
\item
  \textbf{\texttt{w\ =\ (condind\ -\ completeind)\ /\ totalvar}} ---
  within-family variance reduction attributable to ascribed
  characteristics beyond PGIs: (m1 → m2). Hardest detection problem due
  to the large parameter burden (all PGIs, ascribed, and their
  interaction).
\end{enumerate}

\subsubsection{Steps}\label{steps}

\textbf{Step 1 --- Calibrate parameters from observed WLS data} Read
variance components from estimated models.

{\def\LTcaptype{none} % do not increment counter
\begin{longtable}[]{@{}lllll@{}}
\toprule\noalign{}
Sample & N families & ICC & delta\_within & w\_obs \\
\midrule\noalign{}
\endhead
\bottomrule\noalign{}
\endlastfoot
Brothers & 473 & 0.377 & 0.038 & 0.008 \\
Sisters & 612 & 0.373 & 0.033 & 0.025 \\
\end{longtable}
}

\textbf{Step 2 --- Data-generating process} Generates synthetic sibling
datasets with a controlled variance decomposition. Each PGI is split
into a between-family component (f\_j \textasciitilde{} N(0, 0.5)) and a
within-family component (w\_ij \textasciitilde{} N(0, 0.5)). PGI
coefficients are calibrated so that within-family PGI contribution
equals exactly \texttt{delta\_within}. For \texttt{w}, ascribed
variables additionally respect their real-world between/within variance
structure; their coefficients are calibrated so their joint
within-family contribution equals exactly \texttt{w\_true}. Residual
variances \texttt{sigma2\_tau} and \texttt{sigma2\_delta} are set to
reproduce the observed ICC exactly.

\textbf{Step 3 --- Fit models and extract variance components} Fits the
three-model sequence (m0, m1, m2) via \texttt{lmer}, reproducing the
paper's full-interaction specification. Extracts \texttt{delta\_within}
(m0 → m1) and \texttt{w} (m1 → m2).

\textbf{Step 4 --- Null distribution} Runs \texttt{n\_null} replications
with the estimand set to zero (\texttt{delta\_within\ =\ 0} for
\texttt{delta\_within}; \texttt{w\_true\ =\ 0} for \texttt{w}) at the
actual sample size of each subsample. The 95th percentile of the
resulting distribution is the critical value for a one-sided alpha =
0.05 test.

\textbf{Step 5 --- Power curve} Loops over an effect size grid at the
fixed observed N for each sample. For each grid point, runs
\texttt{n\_sim} replications and records the fraction where the estimate
exceeds the null critical value.

\textbf{Step 6 --- Plot} Power curve with x-axis = true effect size,
y-axis = power. A dashed line marks 80\% power; a red dotted line marks
the observed effect size; a blue dotted line marks the minimum
detectable effect (linearly interpolated where the curve crosses 80\%).

\subsubsection{Scripts}\label{scripts}

\begin{itemize}
\tightlist
\item
  \textbf{\texttt{11\_POWER\_d\_simulation.R}} --- power for
  \texttt{delta\_within}; fits m0 and m1 only (no ascribed variables
  required)
\item
  \textbf{\texttt{12\_POWER\_w\_simulation.R}} --- power for \texttt{w};
  fits all three models including the full m2 with \textasciitilde210
  fixed effects
\end{itemize}

\end{document}
